\chapter{TINJAUAN PUSTAKA}
\label{chap:tinjauanpustaka}


\section{Hasil Penelitian Terdahulu}
\label{sec:hasil_terdahulu}

Penelitian oleh \parencite[]{macenski2023desks} membahas implementasi \emph{layered costmap} di dalam library ROS 2, yaitu sebuah metode untuk menyimpan informasi lingkungan secara sistematis dengan tiap lapisan mampu menyimpan informasi tertentu. Tiap lapisan dapat mewakili sumber informasi, algoritma, ataupun hasil berbeda. Standarnya di ROS 2 \emph{layered costmap} dibagi menjadi lapisan static, inflation, obstacle, dan voxel. Manfaatnya dengan menggunakan \emph{layered costmap} termasuk pembaruan yang lebih jelas, area pembaruan yang dinamis, proses pembaruan yang teratur, dan konfigurasi yang fleksibel, memungkinkan representasi nilai cost yang kompleks dan perilaku navigasi yang membutuhkan konteks, seperti navigasi sosial.

Studi oleh \parencite{hu2022slope} mengusulkan pendekatan navigasi adaptif kemiringan untuk robot bergerak darat, yang bertujuan mengatasi keterbatasan costmap 2-dimensi yang sering salah mengartikan kemiringan (slope) sebagai area terlarang. Pendekatan baru ini didasarkan pada layered costmap yang secara aktif mengintegrasikan informasi kemiringan selama tahap pembuatan peta lingkungan. Sebuah algoritma deteksi kemiringan dikembangkan untuk mengidentifikasi kemiringan dan secara adaptif mengubah peta biaya selama navigasi, sehingga kemiringan dianggap sebagai jalur yang dapat dilalui, hanya dengan biaya tambahan. 

Beberapa penelitian sebelumnya mengusulkan metode untuk mewakilkan informasi multi-lantai dari bangunan. Studi oleh \parencite[]{kim2024development} menguraikan pengembangan \emph{indoor delivery mobile robot} dengan arsitektur yang dirancang khusus untuk lingkungan multi-lantai, dengan menggunakan peta navigasi terintegrasi yang dibuat dengan menggabungkan peta \emph{point cloud} multi-lantai dan peta topologi berbasis node graph. Studi oleh \parencite[]{palacin2023procedure} membuat metode navigasi antar lantai di gedung bertingkat dengan menggunakan lift, dengan lokalisasi algoritma Iterative Closest Point (ICP). Performa ICP dievaluasi dalam berbagai kondisi, yaitu saat pintu lift dalam proses terbuka dan tertutup. Sedangkan, \parencite[]{jung2024munes} menggunakan sensor barometer untuk mendeteksi perbedaan ketinggian antar lantai saat menggunakan lift.

Beberapa penelitian membuat metode untuk secara otomatis mendeteksi perantara pergantian lantai seperti lift dan tangga. 
Paper oleh \parencite{li2021autonomous} menunjukan robot melakukan navigasi gedung bertingkat dengan mengoperasikan lift secara otonom. Sedangka, penelitian oleh \parencite[]{lee2023autonomous} dan \parencite[]{kim2020mobile} memanfaatkan pengenalan tombol lift dengan \emph{deep learning} untuk memungkinkan robot beroperasi secara mandiri di antara lantai tanpa memerlukan modifikasi fasilitas lift yang ada atau sistem manajemen gedung tambahan. Penelitian oleh \parencite[]{zhu2020atv} mengusulkan metode \emph{real-time} untuk mendeteksi dan lokalisasi tangga dengan sensor LiDAR VL-16P yang digunakan untuk mendeteksi fitur geometris dari tangga 

Untuk metode \emph{path planning}, \parencite[]{xie2023autonomous} membuat sistem navigasi untuk melintasi lantai dan tangga dengan mengintegrasikan algoritma A* untuk perencanaan jalur awal dan \emph{DWA-based path following} untuk permukaan datar. Paper oleh \parencite[]{jung2024munes} mengusulkan skema untuk perencanaan lintasan menggunakan algoritma A* yang ditambah dengan fungsi biaya untuk mempertimbangkan faktor-faktor realistis seperti waktu tunggu lift, dan ini memungkinkan robot memilih mode pergerakan antar lantai (lift atau tangga). Penelitian oleh \parencite{palacin2023path} menggunakan \emph{static navigation tree} yang menggambarkan \emph{weighted paths} dalam bangunan bertingkat. Lalu, \emph{navigation tree} digunakan untuk merencanakan jalur dan mengestimasi durasi dari semua rute pengiriman menggunakan algoritma dijkstra. 

\section{Teori Dasar}
\label{sec:konsep dasar}

\subsection{Occupancy Grid}

Occupancy Grid adalah salah satu metode pemetaan spasial yang paling fundamental dalam robotika, di mana ruang lingkungan direpresentasikan sebagai sebuah kisi (grid) yang terdiri dari sel-sel berukuran sama. Sejauh ini, domain yang paling umum dari peta \textit{occupancy grid} (kisi keterhunian) adalah peta denah lantai 2-D, yang menggambarkan irisan 2-D dari dunia 3-D. Peta 2-D seringkali sudah cukup, terutama ketika robot bernavigasi di atas permukaan datar dan sensor dipasang sedemikian rupa sehingga hanya menangkap satu irisan dari dunia tersebut. Teknik \textit{occupancy grid} dapat digeneralisasi menjadi representasi 3-D, tetapi membutuhkan beban komputasi yang signifikan.

Misalkan $\mathbf{m}_i$ menyatakan sel kisi dengan indeks $i$. Sebuah peta \textit{occupancy grid} mempartisi ruang menjadi sejumlah sel kisi yang terbatas:

\begin{equation}
m = \sum_i \mathbf{m}_i \tag{9.2}
\end{equation}

Setiap $\mathbf{m}_i$ memiliki nilai keterhunian biner yang terikat padanya, yang menentukan apakah sebuah sel terisi (\textit{occupied}) atau kosong (\textit{free}). Kita akan menuliskan '1' untuk terisi dan '0' untuk kosong \parencite{thrun2002probabilistic}. 

\subsection{Layered Costmap}
Costmap adalah representasi grid dua dimensi dari lingkungan robot yang digunakan dalam sistem navigasi, dengan tiap sel pada grid memiliki harga yang terasosiasi dengannya. Pada \emph{costmap} modern, tiap sel merupakan integer dari 0 hingga 255 yang merepresentasikan biaya navigasi. Secara umum, algoritma perencanaan jalur akan mengutamakan melewati sel dengan harga rendah dan menghindari sel dengan harga tinggi. Metode \emph{costmap} memberikan keseimbangan antara kualitas informasi dan efisiensi algoritma. Sebagai ekstensi dari \emph{costmap}, digunakan konsep \emph{layered costmap}, yaitu struktur costmap yang terbagi menjadi beberapa layer terpisah berdasarkan jenis data atau fungsi spesifik (misalnya static map, obstacle, inflation, proxemic). Setiap layer memodifikasi master costmap secara terstruktur dan independen, memungkinkan sistem navigasi mempertimbangkan konteks yang lebih kompleks dan menghasilkan perilaku pergerakan robot yang lebih cerdas dan adaptif terhadap lingkungan \parencite[]{macenski2023desks}.

\subsection{Lokalisasi Iterative Closest Point (ICP)}
Lokalisasi robot adalah proses untuk menentukan posisi dan orientasi robot (pose) relatif terhadap peta lingkungan tempat robot beroperasi. Kemampuan ini sangat penting bagi robot untuk melakukan navigasi dan pengambilan keputusan yang tepat dalam menjalankan misi otonom. Lokalisasi umumnya menggunakan informasi dari sensor internal, seperti encoder roda dan IMU, dan sensor eksternal, seperti kamera atau LiDAR, untuk memperkirakan posisi robot berdasarkan model gerak dan observasi lingkungan. 
\par
Salah satu kategori metode lokalisasi adalah map-maptching, dimana data sensor dari lingkungan dicoba untuk disamakan dengan fixed map yang sudah ada. Algoritma iterative closest point (ICP) merupakan salah satu algoritma yang bekerja dengan metode map-matching ini. Cara kerja algoritma ICP adalah dengan mencoba mengurangi jarak Euclidean antara input data dengan sebuah model referensi (dalam lokalisasi adalah fixed map). Secara matematis, bayangkan terdapat dua set titik 2D A dan B. Tujuan algoritma adalah menemukan sebuah transformasi A ke B untuk mencari mean squared distance (MSD) antara titik A dan B \parencite[]{sobreira2019map}.
\begin{equation}
\text{MSD}(A, B, u) = \frac{1}{n} \sum_{a \in A, b \in B} \left\| b - u(a) \right\|^2
\end{equation}

Dengan rumus matematika diatas, algoritma ICP mencoba untuk mengurangi MSD pada setiap iterasi. ICP terbagi menjadi dua tahap, yaitu matching stage dan transformation stage. Matching stage adalah tahap pertama dari ICP dan bertujuan untuk meminimalkan nilai mean square distance dengan dengan matching terbaik antara set titik A dan B.
\par
Transformation stage adalah tahap kedua yang bertujuan untuk mencari nilai paling optimal dari transformasi rotasi R dan translasi t untuk mencapai MSD paling kecil. ICP menggunakan simple least square solver untuk menemukan transformasi matrix (R | t) yang meminimalkan MSD. untuk mencapai hal tersebut dilakukan pengurangan dari data referensi dan point cloud sensor nilai centroid masing-masing, yang ditunjukan di rumus dibawah \parencite[]{sobreira2019map}.

\begin{equation}
a'_i = a_i - \frac{1}{n} \sum_{i=0}^{n} a_i
\end{equation}

\begin{equation}
b'_i = b_i - \frac{1}{m} \sum_{i=0}^{m} b_i
\end{equation}

Setelah itu, nilai kovarian silang matrix dikomputasi menggunakan persamaan dibawah dengan $A'$ dan $B'$ sama dengan set poin $a'_i$ dan $b'_i$.
\begin{equation}
H = A'B'^T
\end{equation}
Lalu, sudut rotasi $\theta$ bisa dihitung dengan rumus dibawah .
\begin{equation}
\theta = \text{atan2}((H(0, 1) - H(1, 0)), (H(0, 0) + H(1, 1)))
\end{equation}

Pada akhir tahap transformasi, data sensor diubah menggunakan matriks (R | t) yang diperkirakan dan algoritma kembali ke tahap pencocokan, kecuali jika kriteria penghentian diverifikasi (misalnya, jumlah iterasi, peningkatan kesalahan Euclidean antar iterasi, dll.) \parencite[]{sobreira2019map}.
\begin{equation}
R = m \begin{bmatrix} \cos(\theta) & -\sin(\theta) \\ \sin(\theta) & \cos(\theta) \end{bmatrix}
\end{equation}
\begin{equation}
t = \bar{b} - R\bar{a}
\end{equation}

\subsection{Path Planning A* }
\emph{Path planning} adalah proses perencanaan jalur geometris yang menghubungkan titik awal ke titik tujuan tanpa menabrak rintangan, biasanya dilakukan dalam ruang konfigurasi (configuration space) dengan mempertimbangkan berbagai batasan lingkungan. Metode \emph{path planning} umum mencakup teknik roadmap, cell decomposition, dan artificial potential field.
\par
Salah satu algoritma \emph{path planning} yang paling sering digunakan adalah A*. A* dikategorikan sebagai sebagai algoritma best first search (BFS) yang menggabungkan kelebihan uniform-cost dan greedy searches dengan menggunakan fitness function. 
\begin{equation}
    f(n) = g(n) + h(n)
\end{equation}
\par
Pada persamaan diatas, g(n) adalah biaya kumulatif dari titik awal ke simpul n, dan h(n) adalah estimasi heuristik biaya tersisa menuju simpul tujuan. Selama pencarian, A* mengelola daftar terbuka (simpul yang akan dipertimbangkan) dan daftar tertutup (simpul yang sudah dikunjungi). Algoritma ini bekerja dengan memperluas simpul dari daftar terbuka dengan fungsi kecocokan minimal, memindahkannya ke daftar tertutup, dan menambahkan tetangganya ke daftar terbuka hingga simpul tujuan diperluas.
\par
Pilihan heuristik yang baik sangat penting untuk kualitas dan efisiensi pencarian A*. Heuristik yang mengestimasi biaya nyata menjamin jalur terpendek, tetapi bisa memperluas terlalu banyak simpul. Sebaliknya, heuristik yang melebih-lebihkan dapat mempercepat pencarian ke tujuan, namun berisiko melambatkan jika jalur optimal menyimpang. Meskipun A* banyak digunakan karena menemukan jalur biaya minimum dan memastikan keberadaan jalur bebas, kelemahan utamanya adalah konsumsi waktu yang tinggi dan berat komputasi yang signifikan \parencite[]{damle2022dynamic}.

\subsection{Trajectory Generator}
Pembuatan trajectory atau \emph{Path tracking} bertugas menambahkan informasi waktu pada jalur tersebut, sehingga menghasilkan lintasan yang dapat diikuti oleh robot secara fisik. Proses ini mempertimbangkan batasan kinematik dan dinamik dari robot, termasuk kecepatan, percepatan, dan jerk, untuk menghasilkan lintasan yang halus, aman, dan efisien. \emph{path tracking} memastikan bahwa gerakan aktual robot mengikuti jalur yang telah direncanakan dengan akurat dan memastikan kestabilan robot serta performa yang sesuai dengan sistem kontrol yang digunakan \parencite[]{yao2020control}.

Untuk membuat gerakan robot yang mulus (tidak menyentak atau jerky), kita tidak bisa sekadar memindahkan kecepatan dari 0 ke 1 secara mendadak. Dua metode trajectory yang paling umum digunakan:
\begin{enumerate}
  \item \textbf{Polinomial (Cubic dan Quintic)}: Fungsi waktu dibuat menggunakan persamaan polinomial. Polinomial orde-3 (Cubic) digunakan untuk memastikan kecepatan robot di titik awal dan titik akhir bernilai nol. Polinomial orde-5 (Quintic) selangkah lebih maju dengan memastikan percepatannya juga nol di titik awal dan akhir, sehingga pergerakan jauh lebih halus bagi motor robot.
  \item \textbf{Profil Kecepatan Trapesium}: Robot berakselerasi dengan nilai konstan di awal, melaju dengan kecepatan konstan di pertengahan jalur, lalu melambat dengan deselerasi konstan di akhir. Jika digrafikkan terhadap waktu, kecepatannya akan membentuk bangun trapesium. Ini adalah metode yang paling banyak digunakan di industri.
\end{enumerate}

Dalam praktiknya, robot jarang hanya bergerak dari titik A ke titik B. Robot sering harus melewati banyak titik (waypoints). Jika robot harus benar-benar berhenti di setiap titik, gerakannya akan sangat kaku dan memakan waktu. Daripada berhenti total, lintasan robot dikalkulasi agar bisa "memotong sudut" di dekat waypoint tanpa mengurangi kecepatan hingga nol, menjaga momentum dan fluiditas pergerakan \parencite{lynch2017modern}.

\subsection{ROS}

Robot Operating System (ROS) merupakan sebuah kerangka kerja middleware opensource yang menyediakan navigation stack komprehensif dengan mengintegrasikan fungsi persepsi lingkungan, konstruksi peta, perencanaan jalur, dan kontrol gerakan. Sistem pada ROS beroperasi menggunakan arsitektur modular yang memisahkan antara perangkat keras dan perangkat lunak. Modul-modul fungsional di dalam ROS dihubungkan melalui mekanisme komunikasi publish-subscribe, yang memungkinkan terjadinya pemisahan modul secara fleksibel, penggantian plugin, dan penyesuaian parameter secara langsung saat sistem sedang berjalan. Arsitektur closed-loop ini menyediakan fondasi yang sangat tangguh dan adaptif untuk mendukung operasional navigasi otonom pada robot seluler.

Di dalam arsitektur navigasi ROS, lapisan path planning berperan penting dalam mengarahkan pergerakan robot dengan aman, yang pada umumnya terdiri dari dua komponen utama, perencanaan jalur global (seperti A* dan Dijkstra) serta penghindaran rintangan lokal (seperti DWA dan TEB). Algoritma-algoritma pada lapisan ini secara dinamis menghasilkan lintasan yang optimal berdasarkan pemodelan peta saat ini serta informasi rintangan yang ada di sekitarnya. Di samping itu, lapisan perencanaan ini juga dirancang untuk mendukung kemampuan perencanaan ulang secara waktu nyata, sehingga robot dapat merespons dan menghindari rintangan dinamis secara efektif di dalam lingkungan yang kompleks.

Perencana jalur global (global planner) bertugas untuk mengkalkulasi rute keseluruhan secara makro pada peta statis, di mana algoritma A* dan Dijkstra adalah metode yang paling luas penggunaannya. Untuk mengeksekusi rute global sekaligus menghindari rintangan lokal secara langsung, sistem ROS umumnya memanfaatkan Dynamic Window Approach (DWA) sebagai plugin perencana lokal bawaan karena algoritma ini memiliki beban komputasi yang rendah dan performa waktu nyata yang bagus. DWA menghasilkan lintasan yang aman dan layak secara waktu nyata dengan cara melakukan pengambilan sampel pada ruang kecepatan. Setiap sampel kecepatan hasil tersebut kemudian dievaluasi secara matematis berdasarkan fungsi biaya yang mencakup perhitungan keselarasan arah atau, jarak kedekatan dengan rintangan, dan efisiensi kecepatan.

Dalam proses pelacakan jalur, DWA bekerja dengan memprediksi posisi robot di masa depan berdasarkan masukan kecepatan saat ini. Algoritma ini kemudian mengkalkulasi deviasi arah dan jarak minimum terhadap rintangan terdekat sebagai acuan untuk memandu sistem dalam memilih lintasan terbaik. Setelah lintasan lokal terbaik terpilih, lapisan eksekusi kontrol akan menerjemahkan lintasan tersebut menjadi perintah kendali gerak yang konkret. Perintah berupa kecepatan linier dan sudut ini dipublikasikan secara kontinu melalui antarmuka standar perintah ROS yaitu cmd\_vel, yang pada akhirnya membentuk sistem kendali tertutup dengan base controller robot penggerak \parencite{wei2025ROS}.


